% Chapter 7: Security Analysis

\chapter{Security Analysis}

This chapter analyzes the security properties of PVE-ZKP under the system model
and threat assumptions introduced earlier. It focuses on correctness, soundness
against malicious clients, data minimization, and trust separation, and sketches
how PVE-ZKP can be combined with existing circuit-analysis tools.

\section{Correctness}

Correctness requires that honest executions with valid inputs lead to
acceptance by the receiving system. Assuming correct implementation of the
policy circuits and Groth16, any form submission that satisfies all validation
policies will yield proofs that verify. The receiving system, which bases its
decision solely on the verification outcome, will accept such submissions.

Bugs in circuit code or in the mapping from form fields to circuit inputs could
violate correctness by causing valid submissions to be rejected. These issues
are engineering defects rather than cryptographic failures, and can be mitigated
through testing and code review.

\section{Soundness Against Malicious Clients}

Soundness ensures that a malicious client cannot cause the receiving system to
accept a submission that does not satisfy the specified policies. In PVE-ZKP,
this property relies on the soundness of the underlying Groth16 proof system
and the absence of underconstrained signals in the policy circuits.

If a client could find an assignment to circuit inputs that violates the
intended policy while still satisfying all constraints, or forge a proof without
knowing a valid witness, it might cause the receiving system to accept invalid
inputs. Circuit-analysis tools that search for underconstrained signals and
incomplete constraint sets are therefore valuable complements to the PVE-ZKP
engineering process.

\section{Data Minimization}

PVE-ZKP is designed to minimize the exposure of raw user data. The receiving
system never sees the raw form data submitted by the browser; instead, it sees
only proofs and public signals. This limits the information held by the
backend and separates policy enforcement from data disclosure.

However, the existence of a proof still reveals that some input satisfying the
policy exists, and timing, error messages, or access patterns could leak
additional information. PVE-ZKP does not attempt to provide formal differential
privacy guarantees, and developers must consider application-specific leakage
channels.

\section{Trust Separation}

The architecture separates the roles of validator and receiver. The validator
may be operated by a third party and is not assumed to be fully trusted. Even
if the validator misbehaves or is compromised, the receiving system will only
accept submissions for which it can verify valid proofs.

Trust separation is not absolute: the validator still sees raw inputs and could
log or misuse them. PVE-ZKP reduces---but does not eliminate---the need to
trust the validator, and additional organizational and technical controls are
needed to govern validator behavior in practice.

\section{Integration with Circuit-Analysis Tools}

As noted in the background chapter, recent research has produced tools for
analyzing ZK circuits for security issues. Integrating such tools into the
PVE-ZKP development cycle helps ensure that the policy circuits accurately
encode the intended validation logic and that no unintended degrees of freedom
exist.

A mature deployment of PVE-ZKP would include automated checks that run these
analysis tools on each circuit whenever it is modified, as part of a continuous
integration pipeline.

\section{Residual Risks}

Residual risks include side-channel leakage, misconfiguration of environment
variables and secrets, and subtle bugs in glue code that maps form inputs to
circuit signals. The current prototype does not address these issues fully; they
are discussed further in the next chapter on limitations and future work.
